\chapter{Introduction}\label{ch:introduction}

\section{When code becomes cheap, the specification becomes scarce}

For most of the history of software engineering, source code has been the scarce,
durable, and authoritative artefact of a project. Code was expensive to produce
and to change; it accreted the team's institutional memory---the handling of an
awkward edge case, the workaround for a third-party defect, the carefully tuned
boundary of a loop---and it tended to outlive the informal notes that described
it. The intent behind the code, what the program was \emph{supposed} to do, was
usually recorded only informally and routinely allowed to drift out of step with
the implementation.

The maturation of large language models for code generation disturbs this settled
economics. Systems trained on vast corpora can now produce working implementations
from natural-language descriptions \citep{chen2021codex,li2022alphacode}, and
agentic variants can navigate a repository, edit files, run a build, and iterate
toward a passing test suite with limited human intervention
\citep{yang2024sweagent,jimenez2024swebench}. As the marginal cost of producing a
plausible implementation falls toward zero, the relative value of a project's
artefacts inverts. If code can be regenerated cheaply from a sufficiently precise
description of intent, it becomes a disposable by-product, and the durable, scarce
artefact---the thing that must be preserved, reviewed, and versioned---becomes the
precise \emph{specification} against which code is generated and judged. The claim
is deliberately conditional: \emph{to the extent that} code generation becomes
routine, the bottleneck migrates from writing code to specifying intent precisely
enough that generated code can be trusted. Model output is fluent but unreliable;
it hallucinates interfaces and produces code that passes a weak test suite while
remaining subtly wrong \citep{zhang2025hallucination,liu2023evalplus}. The
discipline's traditional answer to ``how do we know this code is correct?''---read
it carefully---scales poorly when code is produced faster than it can be read. What
is required instead is a machine-checkable statement of what the code must do, and
a mechanical means of checking it. That statement is a formal specification.

\section{The barrier: specifications are valuable but unwritten}

That programs should carry precise, machine-checkable statements of their intended
behaviour is among the oldest aspirations of the field. Hoare's axiomatic basis
\citep{hoare1969axiomatic} and Dijkstra's discipline of programming
\citep{dijkstra1976discipline} established preconditions, postconditions, and loop
invariants as the vocabulary of correctness; Meyer's design by contract
\citep{meyer1992dbc} made the contract, rather than the implementation, the unit
of interface between components---precisely the property that lets an
implementation be regenerated without disturbing its callers, provided the
contract is preserved \citep{mostafa2017behavioral}. The Java Modeling Language
(JML) gave this discipline a tool-supported notation \citep{leavens2006jml}, and
modern verifiers such as OpenJML \citep{cok2014openjml}, discharging obligations
through SMT solvers such as Z3 \citep{demoura2008z3}, can check non-trivial Java
methods against JML contracts by extended static checking without running the
program. The machinery for \emph{checking} programs against contracts is mature;
what has been missing is an economical supply of the contracts.

Despite this lineage, formal specifications remain rare in industrial codebases.
This is an \emph{adoption} problem, not a \emph{value} problem. Few engineers
dispute that a precise contract is useful; many never write one, because authoring
specifications is costly, demands unevenly distributed expertise, and yields
benefits that are diffuse and deferred while the costs are immediate and
concentrated on the author. Faced with this asymmetry, the rational engineer under
schedule pressure writes a unit test, or nothing. The same asymmetry recurs for
whole codebases: legacy systems---most systems that matter---are precisely those
whose behaviour is no longer fully understood \citep{feathers2004legacy}, and the
cost of retrofitting specifications by hand is prohibitive in exactly the cases
where they would be most valuable. This diagnosis points at the leverage point. If
the dominant obstacle is the cost of \emph{authoring} specifications, then lowering
that cost---ideally to near zero, by inferring specifications automatically from
existing code---should change the calculus of adoption. Automatic inference is
itself a mature line of research: Daikon infers likely invariants from executions
\citep{ernst2007daikon}, and Houdini guesses candidate annotations and prunes
those that fail to verify \citep{flanagan2001houdini}. What they leave open is
whether inference can be made cheap and unobtrusive enough to run as a matter of
course, whether its output is genuinely useful downstream, and whether the
resulting specifications can be composed and integrated into how software is built.

\section{Why now: agentic lifecycles without a verification loop}

The case for cheap specifications would be academic were it not for a concurrent
shift in how software is built. A new family of development lifecycles has
emerged---AWS's AI-Driven Development Lifecycle \citep{awsaidlc}, the GitHub Spec
Kit \citep{githubspeckit}, and AWS's Kiro \citep{awskiro}---in which a developer
writes a specification of desired behaviour, an agent plans and implements against
it, and the cycle repeats. These lifecycles at last place the specification at the
centre of the workflow, but they contain a structural gap: the specifications that
drive them are written in \emph{natural language}, which is not machine-checkable.
The verification step, where it exists, falls back on tests synthesised by the same
fallible model against the same ambiguous description, and on human review---the
very bottleneck the lifecycle was meant to relieve. The agentic \emph{capability}
was demonstrated first on benchmarks where a hidden test suite supplies the oracle
for free \citep{jimenez2024swebench}; the industrial lifecycles transplant the
capability but leave the oracle behind. The result is a loop that generates code
quickly but cannot \emph{verify} it: there is no machine-checkable contract layer
and no closed loop in which a candidate is checked against a precise oracle and
rejected when it fails.

The compatibility dimension of this gap deserves emphasis. Modern software is
assembled from libraries, and the contract between a library and its clients is
communicated by a version number under the convention of semantic versioning
\citep{preston2013semver}. This convention is honoured unevenly---breaking changes
are routinely shipped in releases the version number declares compatible
\citep{raemaekers2017semver}---because the contract being versioned is implicit:
there is no machine-checkable statement of what a library guarantees, and so no
mechanical way to decide whether a new release still honours what the old one
promised. An agentic lifecycle that regenerates components against natural-language
descriptions inherits and amplifies this problem, as each regeneration is a
potential silent breaking change. A compositional, machine-checkable specification
is exactly what turns version compatibility from an asserted convention into a
checkable property.

\section{Thesis statement and contributions}

The central claim of this thesis is that \textbf{formal specifications are valuable
but unwritten because authoring them is costly; automatic inference makes them
cheap enough first to be useful, then to be composed to reason about library
compatibility, and finally to serve as the contract layer and verification loop of
AI-native development lifecycles.} The two animating concerns, led with throughout,
are formal-methods \emph{adoption}---the removal of the authoring barrier through
inference---and library and version \emph{compatibility}---the use of
compositional, automatically inferred specifications to reason about whether a
change to a component preserves the behaviour its callers depend upon. Two words
guard the statement. ``Cheap'' does not mean free of judgement: inference removes
the up-front authoring cost by beginning from a verified draft rather than a blank
page. And where the thesis proposes a lifecycle architecture, it is careful to
distinguish what has been \emph{measured} from what has been \emph{designed}.

The programme is organised around three research questions that trace this arc.
\textbf{RQ1 (fidelity and utility):} can JML contracts be inferred automatically
with enough fidelity to be useful, and do they measurably improve a downstream task
such as large-language-model test generation? \textbf{RQ2 (compositional
compatibility):} does a compositional weakest-precondition pass extend single-pass
inference, and does it let behavioural version compatibility be decided as a static
analysis across a call graph? \textbf{RQ3 (lifecycle integration):} can inferred,
machine-checkable contracts serve as the contract layer and verification loop of
AI-native development lifecycles? The three contributions that follow answer RQ1,
RQ2, and RQ3 respectively.

The thesis makes three contributions, each developed and evaluated in its own
chapter:

\begin{enumerate}[leftmargin=1.8em]
  \item \textbf{Inferred specifications are useful, and improve LLM test
    generation} (\cref{ch:inferrer}). The JML-Inferrer recovers sound, verifiable
    contracts from Java source by AST heuristics; a controlled comparison shows
    that supplying these contracts to a large language model measurably improves
    the unit tests it generates, relative to the same model working from code
    alone. This establishes that heuristic, partial, automatically inferred
    specifications are nonetheless good enough to deliver downstream utility.

  \item \textbf{Compositional refinement strictly extends inference and enables
    version-compatibility reasoning} (\cref{ch:compositional}). A second,
    weakest-precondition pass over the call graph recovers well-formed
    preconditions that single-pass heuristic inference cannot, and lets a change to
    a callee's contract be propagated to its callers---turning behavioural version
    compatibility into a static, run-free analysis.

  \item \textbf{Inferred contracts as the verification spine of AI-native
    development} (\cref{ch:integration}). A reference architecture in which a
    machine-checkable contract is the pivot between ambiguous human intent and
    precise machine verification, serving as both the input an agent builds against
    and the oracle its output is checked against, and supplying the closed
    verification loop that current agentic lifecycles lack.
\end{enumerate}

\section{The instrument: the JML-Inferrer}

A single instrument gives the programme its empirical coherence. The JML-Inferrer
is a static analysis tool, implemented in Java~21 over the JavaParser abstract
syntax tree library, that infers JML contracts directly from Java source by
heuristics that pattern-match against the syntax tree. Inferred contracts are
checked by OpenJML's extended static checking \citep{cok2014openjml}, with
obligations discharged by Z3 \citep{demoura2008z3}. A static, syntax-driven
approach requires no test inputs, instrumentation, or execution traces, so it
applies to a codebase exactly as it sits in a repository; and because the
heuristics are deterministic functions of the syntax, the same input yields the
same contracts. The heuristics are unapologetically partial: no claim is made that
a recovered contract is complete, only that it is sound, useful, and verifiable.
The discipline that keeps them honest is verification---a candidate contract that
OpenJML cannot discharge is treated as a defect in the inference to be fixed, not
as an acceptable approximation to be shipped. Throughout, the thesis distinguishes
a \emph{verified} specification, which the checker has discharged, from one that is
merely \emph{well-formed}, which is syntactically valid and scope-correct but has
not been through that check, and it never allows the distinction to blur.

\section{Thesis structure}

\Cref{ch:background} establishes the conceptual and technical foundations---design
by contract and its axiomatic antecedents, JML and OpenJML, the SMT machinery, the
prior art in specification inference, and the emerging landscape of LLM code
generation and agentic development---and situates the work against related
research. \Cref{ch:inferrer} presents the JML-Inferrer in detail and reports the
fidelity and downstream-utility results (the first contribution).
\Cref{ch:compositional} develops the compositional weakest-precondition pass and
its application to version-compatibility reasoning (the second contribution).
\Cref{ch:integration} assembles the validated capabilities into a
verification-centric reference pipeline for AI-native development (the third
contribution). \Cref{ch:discussion} discusses the findings and the threats to
their validity, and \cref{ch:conclusion} summarises the contributions, states the
limitations---including the boundary between verified and merely well-formed
specifications---and sets out directions for future work.
